\subsection{\texorpdfstring{Separable extension of the base field in regular algebras \(^{1}\) or normal}{Separable extension of the base field in regular or normal algebras (1)}}
\footnotetext[1]{In this section, an algebra over a field \(k\) is called regular if it is a regular ring. In particular, every extension of \(k\) is a regular algebra. This notion should not be confused with that of a regular extension, introduced in A, V, p. 135, Definition 2, which will not be used here.}

Lemma 1.--- Let A be a Noetherian ring, the union of an increasing filtered family \((\mathrm{A}_{\alpha})_{\alpha\in\mathbb{I}}\) of Noetherian subrings.

\begin{enumerate}
\def\labelenumi{\alph{enumi})}
\item
  If the rings \(\Lambda_{\alpha}\) are regular and if A is a \(\mathrm{A}_{\alpha}\) -module flat for all \(\alpha \in \mathrm{I}\) the ring A is regular.
\item
  If the rings \(\mathrm{A}_{\alpha}\) are normal (§ 1, no. 8), A is normal.
\end{enumerate}

\begin{enumerate}
\def\labelenumi{\alph{enumi})}
\item
  Let \(\mathfrak{m}\) a maximal ideal of A; for every \(\alpha \in I\) let us denote \(\mathfrak{m}_{\alpha}\) the ideal \(\mathfrak{m} \cap A_{\alpha}\) of \(A_{\alpha}\) . Since A is Noetherian, \(\mathfrak{m}\) is of finite type and there exists an element \(\alpha\) of I such that \(\mathfrak{m} = A\mathfrak{m}_{\alpha}\) so that the A-module \((A_{\alpha}/\mathfrak{m}_{\alpha}) \otimes_{A_{\alpha}} A\) is isomorphic to \(A/m\) Since A is flat over \(A_{\alpha}\) , one has \(\mathrm{dp}_{A}(A/m) \leqslant \mathrm{dp}_{A_{\alpha}}(A_{\alpha}/\mathfrak{m}_{\alpha})\) (A, X, p.~141, lcmc 2). Since the rings \(A_{\alpha}\) are regular, it follows that A is regular \(\S 4\) (no. 2, prop. 4).
\item
  Since the \(A_{\alpha}\) are reduced, A is reduced. Let a and b be elements of A such that b is not a zero divisor and the element a/b of the total ring of fractions of A is integral over \(\Lambda\) . There exists a monic polynomial \(P \in A[X]\) such that \(P(a/b) = 0\) . Let \(\alpha\) an element of I such that the ring \(A_{\alpha}\) contains a, b, and the coefficients of P. Since \(A_{\alpha}\) is normal, there exists \(c \in A_{\alpha}\) such that a = bc. Thus a/b = c belongs to \(\Lambda\) and \(\Lambda\) is normal.
\end{enumerate}

Lemma 2. -- Let \(k\) a field, K and L extensions of \(k\) Suppose that K is of finite type and that one of the extensions K or L is separable. Then the ring \(\mathrm{K} \otimes_{k} \mathrm{L}\) is regular.

Let A denote the ring K ⊗k L; it is Noetherian by the corollary of Prop. 2 (no. 2). Suppose first that the extension K is separable. By A, V, p.~130, corollary, there exists a transcendence basis t = (t₁, \ldots..., tₙ) of K such that K is a finite separable extension of the pure extension k(t). The ring E = k(t) ⊗k L, isomorphic to a ring of fractions of a polynomial ring over L, is regular (§ 4, n° 2, cor. 5 of prop. 4); for every prime ideal p of E, the ring κ(p) ⊗E A is isomorphic to κ(p) ⊗k(t) K, hence to a finite product of fields (A, V, p.~35, def. 1, p.~34, th. 4, and p.~33, prop. 5). Since A is a free E-module, it is a regular ring by the cor. of prop. 9 of § 4, n° 5.

Suppose now that L is separable over k. By the first part of the proof, the ring \(K \otimes_{k} L'\) is regular for every finitely generated subextension \(L'\) of L. We conclude by applying Lemma 1.

PROPOSITION 5.--- Let k be a field, A a k-algebra, and K an extension of k. Suppose that A is essentially of finite type or that the extension K is of finite type.

\begin{enumerate}
\def\labelenumi{\alph{enumi})}
\item
  If the ring A \(_{(K)}\) is regular (resp. normal), the ring A is regular (resp. normal).
\item
  If the ring \(\Lambda\) is regular (resp. normal) and if the extension K of \(k\) is separable, the ring \(\mathrm{A}_{(\mathrm{K})}\) is regular (resp. normal).
\end{enumerate}

The A-module \(A_{(K)}\) is free, hence faithfully flat; assertion a) therefore follows from the corollary of Prop. 8 in I, § 3, No. 5 and from Prop. 8 of § 4, No. 5 (resp. from Cor. 2 of Th. 4 of § 1, No. 10). Under the hypotheses of b), for every prime ideal p of A, the ring \(\kappa(\mathfrak{p})\otimes_{k}K\) is regular (lemma 2), and a fortiori normal (§ 4, no. 2, cor. 2 of prop. 4). Assertion b) therefore follows from the corollary of prop. 9 of § 4, no. 5 (resp. from cor. 3 of th. 4 of § 1, no. 10).
% semantic-fragments: legacy-input-seam
\relax{}
